%% -------------------------------------------------------
%% Kapitel 0: Abstract (DE + EN)
%% -------------------------------------------------------
\chapter*{Kurzfassung}
\addcontentsline{toc}{chapter}{Kurzfassung}

Die vorliegende Diplomarbeit beschreibt die Entwicklung von \textbf{DoggoGO}, einer standortbasierten
Plattform für Hundehalter\-innen und Hundehalter. Das Ziel der Anwendung ist es, relevante
Informationen – wie die Positionen von Freilaufwiesen, Sackerlspendern und Hundekot-Müllbehältern
sowie rassenspezifische gesetzliche Vorschriften (Leinen- und Maulkorbpflicht) – in Echtzeit und
standortgebunden bereitzustellen.

Dieser Teil der Arbeit behandelt die technische Umsetzung des \textbf{Backends} auf Basis von
\textbf{ASP.NET~Core~8} sowie die Entwicklung der \textbf{iOS-Applikation} mit \textbf{Swift} und
\textbf{SwiftUI}. Das Backend stellt eine RESTful~API bereit, die von der Angular-Webapplikation
(von Manuel Rankl) und der iOS-App konsumiert wird. Als Persistenzschicht kommt \textbf{PostgreSQL}
in Kombination mit \textbf{Entity~Framework~Core~9} zum Einsatz. Die gesamte Serverinfrastruktur
wird mittels \textbf{Docker} containerisiert und betrieben.

Die iOS-App ermöglicht die Kartenansicht mit standortbasierten Markierungen über \textbf{MapKit},
die Verwaltung von Hundeprofilen sowie die Anzeige rassenspezifischer Rechtsvorschriften. Die
Authentifizierung erfolgt über \textbf{JSON~Web~Tokens~(JWT)}.

\vspace{1.5em}
\noindent\textit{Schlagwörter: ASP.NET Core, Entity Framework Core, PostgreSQL, Docker, iOS, SwiftUI,
MapKit, JWT, REST~API, Geofencing}

\clearpage

%% -------------------------------------------------------
%% Abstract (English)
%% -------------------------------------------------------
\chapter*{Abstract}
\addcontentsline{toc}{chapter}{Abstract}

This diploma thesis describes the development of \textbf{DoggoGO}, a location-based platform for
dog owners. The application aims to provide relevant information in real time – such as the
locations of off-leash areas, poop-bag dispensers, and dog waste bins, as well as
breed-specific legal regulations (leash and muzzle requirements) – based on the user's current
location.

This part of the thesis covers the technical implementation of the \textbf{backend} using
\textbf{ASP.NET~Core~8} and the development of the \textbf{iOS application} using \textbf{Swift}
and \textbf{SwiftUI}. The backend exposes a RESTful~API consumed by both the Angular web
application (developed by Manuel Rankl) and the iOS app. \textbf{PostgreSQL} together with
\textbf{Entity~Framework~Core~9} serves as the persistence layer. The entire server
infrastructure is containerised and operated using \textbf{Docker}.

The iOS app provides a map view with location-based markers via \textbf{MapKit}, dog profile
management, and the display of breed-specific legal requirements. Authentication is handled
via \textbf{JSON~Web~Tokens~(JWT)}.

\vspace{1.5em}
\noindent\textit{Keywords: ASP.NET Core, Entity Framework Core, PostgreSQL, Docker, iOS, SwiftUI,
MapKit, JWT, REST~API, Geofencing}

\clearpage
